# I.1. Теория составления расписания
\subsection{Теория составления расписания}
В формальной постановке задача составления расписания (timetabling problem) сводится к распределению конечного набора событий — учебных занятий — по конечному набору ресурсов: моментов времени, аудиторий и преподавателей, при соблюдении заданного набора ограничений. Несмотря на кажущуюся простоту формулировки, задача относится к классу NP-трудных: уже при умеренном числе дисциплин, групп и аудиторий пространство возможных расписаний исчисляется величинами порядка десятков и сотен порядков, что делает полный перебор всех вариантов принципиально неприменимым на практике. Этот факт — отправная точка для всего дальнейшего изложения: именно невозможность гарантированно быстро найти оптимальное решение методом перебора и заставляет применять либо точные методы с отсечением пространства поиска (программирование ограничений), либо эвристические методы, не гарантирующие оптимальность, но способные найти приемлемое решение за разумное время (генетические алгоритмы, имитация отжига, поиск с запретами).

Ограничения, накладываемые на расписание, принято делить на две категории.
\begin{itemize}
\item Жёсткие ограничения (hard constraints) — это условия, нарушение которых делает расписание физически нереализуемым: преподаватель не может одновременно вести два занятия, аудитория не может одновременно использоваться двумя группами, число пар в день не может превышать вместимость учебной сетки. Расписание, нарушающее хотя бы одно жёсткое ограничение, считается недопустимым (infeasible) независимо от того, насколько хорошо оно удовлетворяет остальным критериям. 
\item Мягкие ограничения (soft constraints), напротив, описывают желательные, но не обязательные свойства расписания: компактность занятий в течение дня без больших окон, соответствие вместимости аудитории размеру группы, равномерное распределение нагрузки по дням недели. Нарушение мягкого ограничения не делает расписание недопустимым, но снижает его качество. 
\end{itemize}

Эта классификация — жёсткие ограничения как фильтр допустимости, мягкие как критерий качества среди допустимых решений — лежит в основе практически всех современных подходов к автоматизации составления расписаний и используется как сквозная идея во второй главе настоящей работы при описании каждого из реализованных алгоритмов.

Исторически задачи такого типа решались методами целочисленного линейного программирования, позже — методами локального поиска и метаэвристиками, заимствованными из теории эволюционных вычислений и статистической физики. Появление современных решателей задач выполнимости ограничений (constraint satisfaction) и, в частности, гибридных SAT-решателей вернуло интерес к точным методам, способным не просто находить решение, но и доказывать его оптимальность либо отсутствие решения вовсе. Именно это разнообразие — от точных методов до чисто эвристических — определило выбор четырёх алгоритмов, рассматриваемых в данной работе, и стало основанием для их сравнительного анализа.


\subsection{Анализ существующих инструментов}

Существующие программные решения для автоматизации составления расписания можно условно разделить на три группы. Первая — это узкоспециализированные коммерческие и открытые системы вузовского масштаба (например, системы класса ASC Timetables, FET — Free Timetabling Software, и аналогичные отечественные разработки), ориентированные на конечного пользователя — методиста, который задаёт исходные данные через графический интерфейс и получает готовое расписание без необходимости разбираться в алгоритмической стороне вопроса. Такие системы, как правило, используют один конкретный алгоритм (чаще всего метаэвристику — генетический алгоритм либо локальный поиск) и не предоставляют возможности сравнить его с альтернативными подходами на тех же данных, поскольку это не входит в задачи конечного продукта.

Вторая группа — это академические и исследовательские реализации, публикуемые в рамках научных работ по теории расписаний; именно из этого пласта литературы заимствована классификация ограничений на жёсткие и мягкие, рассмотренная в предыдущем разделе, а также большинство идей, реализованных во второй главе настоящей работы (атрибутный табу-лист, критерий Метрополиса, индексное кодирование решения для генетических алгоритмов). Такие работы, как правило, сосредоточены на одном алгоритме и его модификациях, а не на сравнении принципиально разных парадигм.

Третья группа — это общие библиотеки и решатели для задач комбинаторной оптимизации, не привязанные к предметной области расписаний напрямую: библиотека DEAP для эволюционных вычислений на Python, решатель ограничений CP-SAT в составе Google OR-Tools, а также универсальные методы глобальной оптимизации (например, имитация отжига в составе библиотеки SciPy). Эти инструменты избавляют от необходимости реализовывать базовые механизмы алгоритмов с нуля, но требуют адаптации к конкретной предметной области: большинство таких библиотек спроектированы для непрерывных или слабо структурированных задач, и применение их «как есть» к задаче расписания с дискретными, категориальными переменными (номер аудитории, идентификатор преподавателя) требует существенной переработки представления решения и операторов поиска — этот вопрос подробно разбирается в начале каждой подглавы второй главы при описании конкретной реализации.

Общий недостаток существующих решений всех трёх групп — отсутствие единой, воспроизводимой методики сравнения разных алгоритмов на одной и той же постановке задачи. Коммерческие системы скрывают алгоритмическую кухню от пользователя, академические работы сравнивают модификации одного алгоритма между собой, а универсальные библиотеки предоставляют инструменты, но не готовую методологию сравнения. Именно эта пробел и определяет практическую направленность второй главы настоящей работы: реализовать несколько принципиально разных алгоритмов на единой модели данных, с единым критерием оценки качества решения, и провести их сравнение на одинаковых наборах исходных данных.
